\section{August 31, 2026}

\subsection{Bases and coordinates}

\begin{definition}[Basis]
  A collection of vectors $\{v_1,\ldots,v_n\}$ in a vector space $V$ is
  a \emph{basis} for $V$ if it is linearly independent and spans $V$.
\end{definition}

The two conditions in the definition can be combined into a statement
about representing vectors uniquely.

\begin{proposition}
  A collection $\{v_1,\ldots,v_n\}$ is a basis for $V$ if and only if
  every vector $v\in V$ has a unique representation of the form
  \[
    v=a_1v_1+\cdots+a_nv_n,
    \qquad a_1,\ldots,a_n\in\bb{F}.
  \]
  When $\{v_1,\ldots,v_n\}$ is a basis, the scalars
  $a_1,\ldots,a_n$ are called the \emph{coordinates} of $v$ with respect
  to this basis.
\end{proposition}

\begin{proof}
  Suppose first that $\{v_1,\ldots,v_n\}$ is a basis. Since the
  collection spans $V$, every $v\in V$ has a representation
  \[
    v=a_1v_1+\cdots+a_nv_n.
  \]
  To prove uniqueness, suppose that
  \[
    v=\widetilde a_1v_1+\cdots+\widetilde a_nv_n
  \]
  is another such representation. Subtracting the two equations gives
  \[
    (a_1-\widetilde a_1)v_1+\cdots
      +(a_n-\widetilde a_n)v_n=0.
  \]
  Because $\{v_1,\ldots,v_n\}$ is linearly independent,
  \[
    a_i-\widetilde a_i=0
  \]
  for each $i$. Thus $a_i=\widetilde a_i$ for every $i$, and the
  representation is unique.

  Conversely, suppose every $v\in V$ has a unique representation as a
  linear combination of $v_1,\ldots,v_n$. The existence of such a
  representation for every $v$ shows that the collection spans $V$.
  Now suppose
  \[
    a_1v_1+\cdots+a_nv_n=0.
  \]
  The zero vector also has the representation
  \[
    0=0v_1+\cdots+0v_n.
  \]
  Uniqueness therefore gives $a_1=\cdots=a_n=0$. Hence the collection is
  linearly independent, so it is a basis for $V$.
\end{proof}

\begin{example}[The standard basis of $\bb{F}^n$]
  For $i=1,\ldots,n$, let $e_i\in\bb{F}^n$ be the vector whose $i$th
  coordinate is $1$ and whose other coordinates are $0$. Thus
  \[
    e_1=\begin{bmatrix}1\\0\\\vdots\\0\end{bmatrix},
    \qquad
    e_2=\begin{bmatrix}0\\1\\\vdots\\0\end{bmatrix},
    \qquad\ldots\qquad
    e_n=\begin{bmatrix}0\\0\\\vdots\\1\end{bmatrix}.
  \]
  Every vector in $\bb{F}^n$ has the unique representation
  \[
    \begin{bmatrix}x_1\\x_2\\\vdots\\x_n\end{bmatrix}
    =x_1e_1+x_2e_2+\cdots+x_ne_n.
  \]
  Consequently, $\{e_1,\ldots,e_n\}$ is a basis for $\bb{F}^n$, called
  its \emph{standard basis}.
\end{example}

\begin{example}[The standard basis of $\mathcal P_n(\bb{F})$]
  Recall that $\mathcal P_n(\bb{F})$ is the vector space of polynomials
  over $\bb{F}$ of degree at most $n$. Each polynomial has a unique
  expression
  \[
    p(t)=a_0+a_1t+\cdots+a_nt^n.
  \]
  Hence
  \[
    \{1,t,t^2,\ldots,t^n\}
  \]
  is a basis for $\mathcal P_n(\bb{F})$, called its standard basis.
\end{example}

\begin{proposition}
  Every finite nonempty spanning set for a vector space contains a basis.
\end{proposition}

\begin{proof}
  Let $\{v_1,\ldots,v_p\}$ span $V$. If this collection is linearly
  independent, then it is already a basis. Otherwise, by
  Proposition~\ref{prop:dependence-redundancy}, some $v_k$ is a linear
  combination of the other vectors. It follows that removing $v_k$ does
  not change the span:
  \[
    \operatorname{span}\{v_1,\ldots,v_p\}
    =\operatorname{span}\{v_1,\ldots,v_{k-1},v_{k+1},\ldots,v_p\}.
  \]
  If the remaining collection is linearly independent, it is a basis.
  If not, repeat the same argument. The original spanning set is finite,
  so this process must terminate. The final collection still spans $V$
  and is linearly independent, and is therefore a basis contained in the
  original spanning set.
\end{proof}

\subsection{Linear transformations}

A \emph{transformation}, also called a map or function, from a set $X$
to a set $Y$ is a rule $T\colon X\to Y$ that assigns to every $x\in X$
an element $T(x)\in Y$.

\begin{definition}[Linear transformation]
  Let $V$ and $W$ be vector spaces over the same field $\bb{F}$. A
  transformation $T\colon V\to W$ is \emph{linear} if
  \begin{enumerate}[label=(\roman*)]
    \item $T(u+v)=T(u)+T(v)$ for all $u,v\in V$;
    \item $T(\alpha u)=\alpha T(u)$ for all $u\in V$ and
      $\alpha\in\bb{F}$.
  \end{enumerate}
\end{definition}

Equivalently, $T\colon V\to W$ is linear if
\[
  T(\alpha u+\beta v)=\alpha T(u)+\beta T(v)
\]
for all $u,v\in V$ and $\alpha,\beta\in\bb{F}$.

\begin{example}[Reflection]
  Define $T\colon\bb{R}^2\to\bb{R}^2$ by
  \[
    T\begin{bmatrix}x_1\\x_2\end{bmatrix}
    =\begin{bmatrix}x_1\\-x_2\end{bmatrix}.
  \]
  Geometrically, $T$ reflects each point across the first coordinate
  axis. To check additivity, observe that
  \begin{align*}
    T\left(
      \begin{bmatrix}x_1\\x_2\end{bmatrix}
      +\begin{bmatrix}y_1\\y_2\end{bmatrix}
    \right)
    &=T\begin{bmatrix}x_1+y_1\\x_2+y_2\end{bmatrix} \\
    &=\begin{bmatrix}x_1+y_1\\-x_2-y_2\end{bmatrix} \\
    &=\begin{bmatrix}x_1\\-x_2\end{bmatrix}
      +\begin{bmatrix}y_1\\-y_2\end{bmatrix} \\
    &=T\begin{bmatrix}x_1\\x_2\end{bmatrix}
      +T\begin{bmatrix}y_1\\y_2\end{bmatrix}.
  \end{align*}
  For $\alpha\in\bb{R}$,
  \begin{align*}
    T\left(\alpha\begin{bmatrix}x_1\\x_2\end{bmatrix}\right)
    &=T\begin{bmatrix}\alpha x_1\\\alpha x_2\end{bmatrix} \\
    &=\begin{bmatrix}\alpha x_1\\-\alpha x_2\end{bmatrix} \\
    &=\alpha\begin{bmatrix}x_1\\-x_2\end{bmatrix}
      =\alpha T\begin{bmatrix}x_1\\x_2\end{bmatrix}.
  \end{align*}
  Thus $T$ is linear.
\end{example}
